% Part: lambda-calculus
% Chapter: church-rosser
% Section: beta-eta-reduction

\documentclass[../../../include/open-logic-section]{subfiles}

\begin{document}

\olfileid[fr]{lam}{cr}{be}
\olsection{\texorpdfstring{$\beta\eta$}{βη}-réduction}

La propriété de Church--Rosser vaut pour
la réduction~$\beta\eta$ ($\bered$).

Dans cette section, $M\beredone M'$ désigne
une contraction en un pas :
$M\bredone M'$ ou $M\eredone M'$.

\begin{lem}\ollabel{lem:one-par}
Si $M \beredone M'$, alors
$M \beredpar M'$.
\end{lem}

\begin{proof}
Raisonnons sur la dérivation de
$M\beredone M'$. Une contraction~$\beta$
à la racine est un cas de la quatrième
règle de $\beredpar$, par le
\olref[pbe]{thm:refl}; ses cas de contexte
suivent des règles d'abstraction et
d'application. De même, une contraction~$\eta$
à la racine a la forme
$\lambd[x][Nx]\eredone N$, avec
$x\notin\FV{N}$, conformément à la
\olref[syn][eta]{defn:beredone}.
La réflexivité $N\beredpar N$ et la règle
\olref[pbe]{defn:beredpar5} donnent
$\lambd[x][Nx]\beredpar N$.
Les contractions~$\eta$ dans une
abstraction ou une application se
traitent par induction et par les règles
de congruence, comme dans le
\olref[b]{lem:one-par}.
\end{proof}

\begin{lem}\ollabel{lem:par-red}
Si $M \beredpar M'$, alors
$M \bered M'$.
\end{lem}

\begin{proof}
Procédons par induction sur la dérivation
de $M\beredpar M'$. Pour les quatre
premières règles, la construction du
\olref[b]{lem:par-red} s'applique :
les pas~$\beta$ et les réductions dans
les sous-termes sont inclus dans~$\bered$.

Si la dernière étape applique la règle
\olref[pbe]{defn:beredpar5}, alors
$M=\lambd[x][Nx]$ et $M'=N'$, avec
$x\notin\FV{N}$ et $N\beredpar N'$.
La contraction~$\eta$ donne
$\lambd[x][Nx]\eredone N$;
l'hypothèse d'induction donne
$N\bered N'$. Ainsi,
$M\bered N'$, par transitivité.
\end{proof}

\begin{lem}\ollabel{lem:str}
$\bered$ est la plus petite relation
transitive qui contient $\beredpar$.
\end{lem}

\begin{proof}
On reprend les deux inclusions du
\olref[b]{lem:str}, en utilisant ici
les deux lemmes précédents. Le cas
d'une réduction de longueur zéro est
couvert par la réflexivité de $\beredpar$
énoncée dans le \olref[pbe]{thm:refl}.
\end{proof}

\begin{thm}\ollabel{thm:cr}
$\bered$ vérifie la propriété de
Church--Rosser.
\end{thm}

\begin{proof}
D'après le \olref[pbe]{thm:cr}, la relation $\beredpar$
vérifie la propriété de Church--Rosser. Le \olref{lem:str}
identifie $\bered$ à sa clôture transitive. Le
\olref[dap]{thm:str} donne alors la propriété pour $\bered$.
\end{proof}

\end{document}
